\section{Theoretical Origin of the SBC/F3 Infrared Sector}

\subsection{Finite Observational Projection and Emergent Memory}

A central premise of the SBC framework is that the observable universe does not necessarily coincide with the total causal structure of physical reality.

Let \(\Psi_{\rm tot}(t)\) denote the state of the full causal structure, and let the observable cosmological state be defined as a projection,
\begin{equation}
\Psi_{\rm obs}(t)
=
\mathcal{P}\Psi_{\rm tot}(t),
\end{equation}
where \(\mathcal{P}\) is an observational projection operator associated with a finite observational window.

The complementary, unobserved sector is
\begin{equation}
\Psi_{\rm hid}(t)
=
(1-\mathcal{P})\Psi_{\rm tot}(t).
\end{equation}

If the total state evolves locally in the full causal space,
\begin{equation}
\frac{d\Psi_{\rm tot}}{dt}
=
\mathcal{L}\Psi_{\rm tot},
\end{equation}
then projection onto the observable sector yields an effective equation of the form
\begin{equation}
\frac{d\Psi_{\rm obs}}{dt}
=
\mathcal{P}\mathcal{L}\mathcal{P}\Psi_{\rm tot}
+
\mathcal{P}\mathcal{L}(1-\mathcal{P})\Psi_{\rm tot}.
\end{equation}

The second term represents the influence of the hidden causal sector on the observable dynamics. Since \(\Psi_{\rm hid}\) is not directly resolved within the observational window, its effect appears as a nonlocal memory contribution.

Formally, this can be written as
\begin{equation}
\frac{d\Psi_{\rm obs}}{dt}
=
\mathcal{L}_{\rm eff}\Psi_{\rm obs}(t)
+
\int_{-\infty}^{t}
K(t,t')\Psi_{\rm obs}(t')\,dt',
\end{equation}
where \(K(t,t')\) is an effective memory kernel induced by the eliminated hidden sector.

Thus, in the SBC framework, memory is not introduced as an independent assumption. It emerges from the finite observational projection:
\begin{equation}
{\rm finite\ observational\ window}
\quad
\Longrightarrow
\quad
{\rm incomplete\ state\ description}
\quad
\Longrightarrow
\quad
{\rm effective\ memory}.
\end{equation}

When the hidden-sector influence is stable and progressively screened, the leading memory kernel is expected to take an exponentially damped form,
\begin{equation}
K(t,t')
\propto
\exp[-\lambda(t-t')]\Theta(t-t').
\end{equation}

This provides the theoretical bridge between the foundational SBC postulate of finite observational access and the phenomenological SBC/F3 infrared response.

\subsection{Observable and Hidden Causal Sectors}

The SBC framework assumes that the total causal structure can be decomposed into an observable sector and a hidden sector,
\begin{equation}
\Psi_{\rm tot}
=
\Psi_{\rm obs}
+
\Psi_{\rm hid}.
\end{equation}

The observable component is obtained through a projection operator,
\begin{equation}
\Psi_{\rm obs}
=
\mathcal P \Psi_{\rm tot},
\end{equation}
while the hidden component is
\begin{equation}
\Psi_{\rm hid}
=
(1-\mathcal P)\Psi_{\rm tot}.
\end{equation}

The projection operator satisfies
\begin{equation}
\mathcal P^2=\mathcal P,
\end{equation}
which defines a consistent separation between accessible and inaccessible causal sectors.

The total state is assumed to evolve according to
\begin{equation}
\frac{d\Psi_{\rm tot}}{dt}
=
\mathcal L \Psi_{\rm tot},
\end{equation}
where \(\mathcal L\) is the generator of the underlying causal evolution.

Projecting the total dynamics onto the observable and hidden sectors yields
\begin{equation}
\frac{d\Psi_{\rm obs}}{dt}
=
\mathcal L_{oo}\Psi_{\rm obs}
+
\mathcal L_{oh}\Psi_{\rm hid},
\end{equation}
and
\begin{equation}
\frac{d\Psi_{\rm hid}}{dt}
=
\mathcal L_{ho}\Psi_{\rm obs}
+
\mathcal L_{hh}\Psi_{\rm hid}.
\end{equation}

Here,
\begin{equation}
\mathcal L_{oo}
=
\mathcal P\mathcal L\mathcal P,
\end{equation}
\begin{equation}
\mathcal L_{oh}
=
\mathcal P\mathcal L(1-\mathcal P),
\end{equation}
\begin{equation}
\mathcal L_{ho}
=
(1-\mathcal P)\mathcal L\mathcal P,
\end{equation}
and
\begin{equation}
\mathcal L_{hh}
=
(1-\mathcal P)\mathcal L(1-\mathcal P).
\end{equation}

The hidden sector acts as a reservoir of causal information that is not directly resolved within the observable window.

Solving formally for \(\Psi_{\rm hid}\) and substituting it back into the observable-sector equation generates an effective nonlocal equation,
\begin{equation}
\frac{d\Psi_{\rm obs}}{dt}
=
\mathcal L_{\rm eff}\Psi_{\rm obs}
+
\int_{-\infty}^{t}
K(t,t')
\Psi_{\rm obs}(t')
\,dt'.
\end{equation}

The memory kernel \(K(t,t')\) is therefore not introduced as an independent postulate. It emerges from eliminating the hidden causal sector.

This gives the fundamental SBC hierarchy:
\begin{equation}
{\rm total\ causal\ structure}
\rightarrow
{\rm causal\ projection}
\rightarrow
{\rm hidden\ sector}
\rightarrow
{\rm memory\ kernel}
\rightarrow
{\rm infrared\ phenomenology}.
\end{equation}

The SBC/F3 sector can therefore be interpreted as the cosmological infrared limit of an observable-hidden causal decomposition.

\subsection{Toward a Fundamental SBC Memory Equation}

The observationally preferred SBC/F3 sector suggests that the effective infrared deformation may arise from a deeper memory dynamics.

Let \(\Psi(t)\) denote the effective observable cosmological state within a finite observational window. In the SBC framework, the observed state is not assumed to be fully local in time, but may depend on a causal memory functional,
\begin{equation}
\mathcal{M}(t)
=
\int_{-\infty}^{t}
K(t,t')\,\Psi(t')\,dt',
\end{equation}
where \(K(t,t')\) is a causal memory kernel.

Causality requires
\begin{equation}
K(t,t')=0
\qquad
{\rm for}
\qquad
t'>t .
\end{equation}

A stable memory sector must also satisfy finite normalization,
\begin{equation}
\int_{-\infty}^{t}
|K(t,t')|\,dt'
<\infty ,
\end{equation}
ensuring that remote past states do not generate divergent contributions.

The simplest kernel satisfying causality, normalization, and monotonic memory loss is
\begin{equation}
K(t,t')
=
K_0
\exp[-\lambda(t-t')]
\Theta(t-t'),
\end{equation}
where \(\Theta\) is the Heaviside causal step function.

This kernel obeys the first-order memory-decay equation
\begin{equation}
\frac{\partial K}{\partial t}
=
-\lambda K .
\end{equation}

Consequently, the effective memory amplitude satisfies
\begin{equation}
\frac{d\mathcal{M}}{dt}
=
-\lambda \mathcal{M}
+
\mathcal{S}(t),
\end{equation}
where \(\mathcal{S}(t)\) is a source term associated with the instantaneous perturbative state.

In the infrared regime, where \(\mathcal{S}(t)\) varies slowly compared with the memory-decay scale, the homogeneous solution dominates the leading screened response:
\begin{equation}
\mathcal{M}(t)
\propto
e^{-\lambda t}.
\end{equation}

Mapping the causal depth to a cosmological redshift variable gives an effective response of the form
\begin{equation}
F_{\rm SBC}(z)
\simeq
-\epsilon
\exp[-k_0(1+z)].
\end{equation}

This corresponds to the leading-order SBC/F3 phenomenological branch with
\begin{equation}
n\simeq 1.
\end{equation}

Thus, the empirical preference for \(n\simeq 1\) can be interpreted as the observational imprint of a first-order causal memory kernel.

The more general phenomenological form
\begin{equation}
F_{\rm SBC}(z)
=
-\epsilon
\exp\left[-(k_0(1+z))^n\right]
\end{equation}
may then be understood as a generalized memory-decay law, where deviations from \(n=1\) encode departures from the minimal first-order kernel.

The observational result \(n\simeq1\) therefore selects the simplest stable SBC memory regime.

\subsection{Kernel Selection by Stability and Causal Normalization}

The emergence of an exponentially screened memory kernel can be motivated from minimal physical requirements.

Let the memory kernel depend on the causal separation
\begin{equation}
\Delta t = t-t',
\qquad
\Delta t \geq 0 .
\end{equation}

Causality requires
\begin{equation}
K(\Delta t)=0
\qquad
{\rm for}
\qquad
\Delta t<0 .
\end{equation}

Finite memory requires the normalization condition
\begin{equation}
\int_0^\infty K(\Delta t)\,d\Delta t < \infty .
\end{equation}

Stability requires that the influence of remote past states decreases monotonically,
\begin{equation}
\frac{dK}{d\Delta t}<0 .
\end{equation}

The simplest first-order stable decay law satisfying these conditions is
\begin{equation}
\frac{dK}{d\Delta t}
=
-\lambda K,
\qquad
\lambda>0 .
\end{equation}

Its solution is
\begin{equation}
K(\Delta t)
=
K_0 e^{-\lambda \Delta t}.
\end{equation}

Including causal support explicitly gives
\begin{equation}
K(t,t')
=
K_0 e^{-\lambda(t-t')}
\Theta(t-t') .
\end{equation}

Thus, the exponential kernel is not an arbitrary choice. It is the minimal causal, normalizable and stable memory response.

When mapped to cosmological redshift or causal-depth variables, this kernel generates the leading SBC/F3 infrared correction
\begin{equation}
F_{\rm SBC}(z)
\simeq
-\epsilon e^{-k_0(1+z)} .
\end{equation}

This corresponds to the phenomenological SBC/F3 form
\begin{equation}
F_{\rm SBC}(z)
=
-\epsilon
\exp\left[-(k_0(1+z))^n\right]
\end{equation}
with
\begin{equation}
n\simeq 1 .
\end{equation}

Therefore, the observational preference for \(n\simeq1\) can be interpreted as the signature of the minimal stable causal-memory regime.

\subsection{Exponential Memory Decay and the Origin of the IR-Safe Sector}

The repeated emergence of \(n\simeq 1\) in the observationally viable SBC/F3 region suggests that the infrared correction may originate from a first-order memory-decay process.

Let \(M(\tau)\) denote an effective cosmological memory amplitude, where \(\tau\) parametrizes the causal depth or accumulated observational scale. The simplest stable decay law is

\begin{equation}
\frac{dM}{d\tau}=-\lambda M,
\end{equation}

with solution

\begin{equation}
M(\tau)=M_0 e^{-\lambda \tau}.
\end{equation}

If the relevant infrared causal depth scales approximately as

\begin{equation}
\tau \propto (1+z),
\end{equation}

then the effective SBC/F3 response becomes

\begin{equation}
F_{\rm SBC}(z)
\simeq
-\epsilon \exp[-k_0(1+z)].
\end{equation}

This is the leading-order case of the phenomenological form

\begin{equation}
F_{\rm SBC}(z)
=
-\epsilon
\exp\left[-\left(k_0(1+z)\right)^n\right],
\end{equation}

with

\begin{equation}
n\simeq 1.
\end{equation}

Thus, the observational preference for \(n\simeq 1\) can be interpreted as the signature of a first-order exponentially screened memory response.

In this interpretation, \(\epsilon\) measures the amplitude of the effective memory imprint, while \(k_0\) controls the infrared screening scale. The preferred region,

\begin{equation}
\epsilon\simeq 0.06,
\qquad
k_0\simeq 0.015-0.022,
\qquad
n\simeq 1,
\end{equation}

corresponds to a weak, slowly decaying, large-scale memory sector.

This provides a physical rationale for why the SBC/F3 deformation is observationally viable: it is sufficiently weak to preserve the \(\Lambda\)CDM background and high-\(\ell\) acoustic structure, while remaining large enough to generate controlled infrared perturbative effects.
